\documentclass[10pt]{article}

\usepackage[utf8]{inputenc}
\usepackage[T1]{fontenc}
\usepackage{amsmath,amssymb}
\usepackage{booktabs}
\usepackage{enumitem}
\usepackage[margin=0.78in]{geometry}
\usepackage[most]{tcolorbox}
\usepackage{xcolor}
\usepackage{hyperref}

\definecolor{ResultBlue}{RGB}{27,79,114}
\hypersetup{colorlinks=true,linkcolor=blue,urlcolor=blue}
\setlist{itemsep=2pt,topsep=3pt,parsep=1pt,partopsep=0pt}
\setlength{\parindent}{0pt}
\setlength{\parskip}{4pt}
\newcommand{\R}{\mathbb{R}}
\newcommand{\dd}{\,\mathrm{d}}
\newcommand{\norm}[1]{\left\lVert#1\right\rVert}
\newcommand{\solution}[1]{\begin{tcolorbox}[breakable,colback=ResultBlue!3,colframe=ResultBlue,
  boxrule=0.8pt,arc=1mm,title=\textbf{Result},fonttitle=\bfseries]#1\end{tcolorbox}}

\title{Summer Bootcamp -- Session 2\\Corrections to the Calculus and ODE Exercises}
\author{Nathan Sauldubois}
\date{August 2026}

\begin{document}
\maketitle

\section*{Correction Guidelines}
The solutions follow the numbering of the exercise sheet. Numerical values are
rounded only at the end. For programming questions, the displayed formulas are
the reference implementation; a vectorized implementation is equally valid.

\section{Python and Elementary Financial Calculations}

\paragraph{Exercise 1 -- Coupon-bond yield, duration, and repricing.}
Differentiation gives
\[
B'(y)=-4e^{-y}-8e^{-2y}-312e^{-3y},\qquad
B''(y)=4e^{-y}+16e^{-2y}+936e^{-3y}.
\]
Since $B'(y)<0$, the equation $B(y)=96.50$ has at most one solution. Moreover,
$B(0)=112$ and $B(0.2)\simeq63.03$, so $[0,0.2]$ brackets it. Bisection or
Newton's iteration
\[
y_{n+1}=y_n-\frac{B(y_n)-96.50}{B'(y_n)}
\]
gives
\solution{$y=0.0515698674$, i.e. a continuously compounded yield of
$5.15699\%$. At this yield,
\[
D_{\rm mod}=2.88387666,\qquad \mathcal C=8.49811801.
\]}
For a yield move $\Delta y$, duration and duration--convexity give
\[
B_{D}=B(1-D_{\rm mod}\Delta y),\qquad
B_{DC}=B\left(1-D_{\rm mod}\Delta y+\tfrac12\mathcal C(\Delta y)^2\right).
\]
\begin{center}
\begin{tabular}{rrrr}
\toprule
$\Delta y$ & Exact price & Duration & Duration--convexity\\
\midrule
$-100$ bp & 99.3244 & 99.2829 & 99.3239\\
$-25$ bp  & 97.1983 & 97.1957 & 97.1983\\
$+25$ bp  & 95.8068 & 95.8043 & 95.8068\\
$+100$ bp & 93.7577 & 93.7171 & 93.7581\\
\bottomrule
\end{tabular}
\end{center}
The quadratic correction is visibly more accurate, especially for $100$ bp.

\section{Explicit Derivatives}

\paragraph{Exercise 2 -- One-dimensional derivatives.}
The derivatives and domains are
\[
f_1'=6x\cos(3x^2-1),\quad \operatorname{Dom}(f_1)=\R,
\]
\[
f_2'=e^{-3x}(2x-3x^2),\qquad
f_2''=e^{-3x}(2-12x+9x^2),\quad \operatorname{Dom}(f_2)=\R,
\]
\[
f_3'=\frac{(1+x^2)/x-2x\log x}{(1+x^2)^2},
\quad \operatorname{Dom}(f_3)=(0,\infty),
\]
\[
f_4'=\frac{e^x}{1+e^x},\qquad
f_4''=\frac{e^x}{(1+e^x)^2},\quad \operatorname{Dom}(f_4)=\R,
\]
\solution{$\displaystyle f_5'=\frac{x}{\sqrt{1+x^2}}$ and
$\operatorname{Dom}(f_5)=\R$.}

\paragraph{Exercise 3 -- Gradient, Hessian, and Jacobian.}
One obtains
\[
\nabla f=\binom{4x-y+4}{-x+6y},\qquad
H_f=\begin{pmatrix}4&-1\\-1&6\end{pmatrix}.
\]
Furthermore,
\[
J_g(x,y)=\begin{pmatrix}e^x\cos y&-e^x\sin y\\1&-2y\end{pmatrix}.
\]

\paragraph{Exercise 4 -- Multivariate chain rule by two methods.}
Here
\[
F(x,y)=(x+y^2)e^{xe^y}.
\]
Direct differentiation yields
\[
\partial_xF=e^{xe^y}\bigl[1+(x+y^2)e^y\bigr],
\]
\[
\partial_yF=e^{xe^y}\bigl[2y+x(x+y^2)e^y\bigr].
\]
On the other hand,
\[
J_\Phi=\begin{pmatrix}1&2y\\e^y&xe^y\end{pmatrix},\qquad
\nabla q(u,v)=\binom{e^v}{ue^v}.
\]
Thus $\nabla F=J_\Phi^\top\nabla q(\Phi)$ reproduces the two formulas above.
Equivalently,
\[
D\Phi h=\binom{h_1+2yh_2}{e^yh_1+xe^yh_2},\qquad
Dq(u,v)k=e^vk_1+ue^vk_2,
\]
and substituting $(u,v)=\Phi(x,y)$ gives the same coefficients of $h_1,h_2$.

\section{Explicit Integration}

\paragraph{Exercise 5 -- One-dimensional integrals.}
Direct integration, substitution, integration by parts, and a trigonometric
substitution respectively give
\solution{\[
I_1=1,\qquad
I_2=\begin{cases}\dfrac{(1+T)^{1-r}-1}{1-r},&r\ne1,\\[4pt]
\log(1+T),&r=1,\end{cases}
\qquad I_3=\frac{e-1}{2},
\]
\[
I_4=-\frac14,\qquad I_5=\frac{\pi}{4}.
\]}
For $I_2$, the limit as $r\to1$ follows from
$\lim_{a\to0}((1+T)^a-1)/a=\log(1+T)$. For $I_4$, the boundary term
$x^2\log x/2$ tends to zero at the origin.

\paragraph{Exercise 6 -- The Gaussian integral.}
Tonelli's theorem allows the nonnegative integrand to be written as
\[
G^2=\iint_{\R^2}e^{-(x^2+y^2)}\dd x\dd y.
\]
With $x=r\cos\theta$, $y=r\sin\theta$ and Jacobian $r$,
\[
G^2=\int_0^{2\pi}\int_0^\infty e^{-r^2}r\dd r\dd\theta
=2\pi\left[-\frac12e^{-r^2}\right]_0^\infty=\pi.
\]
Since $G>0$,
\solution{$\displaystyle G=\sqrt\pi$, and, after $x=\sqrt2\sigma u$,
\[
\int_{-\infty}^{\infty}e^{-x^2/(2\sigma^2)}\dd x
=\sigma\sqrt{2\pi}.
\]
Therefore the Gaussian normalizing constant is
$1/(\sigma\sqrt{2\pi})$.}

\paragraph{Exercise 7 -- Multiple integrals and change of variables.}
For $J_1$, direct integration gives
\[
J_1=\int_{-1}^2\left(6x-\frac92\right)\dd x=-\frac92.
\]
Reversing the order, $0\le y\le3$ and $-1\le x\le2$, gives the same value.
For the triangle $D$,
\[
D=\{0\le x\le2,\ 0\le y\le(2-x)/2\}
=\{0\le y\le1,\ 0\le x\le2-2y\},
\]
and either description gives $J_2=4/3$. Finally, polar coordinates give
\solution{\[
J_3(R)=\int_0^{2\pi}\int_1^Rr^4r\dd r\dd\theta
=\frac{\pi}{3}(R^6-1),\qquad J_3(3)=\frac{728\pi}{3}.
\]}

\section{Black--Scholes Greeks}

\paragraph{Exercise 8 -- Derive the Greeks.}
Recall that
\[
\Phi(x)=\frac{1}{\sqrt{2\pi}}\int_{-\infty}^x e^{-u^2/2}\dd u,
\qquad
\phi(x)=\Phi'(x)=\frac{1}{\sqrt{2\pi}}e^{-x^2/2}.
\]
With all unmentioned variables held fixed, the Greeks are defined by
\[
\Delta=\frac{\partial C}{\partial S},\qquad
\Gamma=\frac{\partial^2 C}{\partial S^2},\qquad
\mathrm{Vega}=\frac{\partial C}{\partial\sigma},
\]
\[
\Theta=-\frac{\partial C}{\partial T},\qquad
\mathrm{Rho}=\frac{\partial C}{\partial r}.
\]
Here $T$ is time to maturity, so the market convention for calendar-time
Theta is $-C_T$.

\textbf{Question 2 -- Key identity.} Since $d_1-d_2=\sigma\sqrt T$ and
\[
d_1+d_2=\frac{2(\log(S/K)+rT)}{\sigma\sqrt T},
\]
we have
\[
\frac{\phi(d_2)}{\phi(d_1)}
=\exp\left(\frac{d_1^2-d_2^2}{2}\right)
=\exp\bigl(\log(S/K)+rT\bigr).
\]
Therefore
\solution{$\displaystyle S\phi(d_1)=Ke^{-rT}\phi(d_2)$.}

\textbf{Question 3 -- Closed-form Greeks.}

\emph{Delta and Gamma.} We use
\[
\partial_Sd_1=\partial_Sd_2=\frac{1}{S\sigma\sqrt T}.
\]
The chain rule gives
\[
\frac{\partial C}{\partial S}
=\Phi(d_1)+\frac{\phi(d_1)}{\sigma\sqrt T}
-\frac{Ke^{-rT}\phi(d_2)}{S\sigma\sqrt T}.
\]
The last two terms cancel by the key identity. Differentiating once more,
\[
\boxed{\Delta=\Phi(d_1),\qquad
\Gamma=\frac{\phi(d_1)}{S\sigma\sqrt T}.}
\]

\emph{Vega.} Differentiating both normal-distribution terms with respect to
$\sigma$ and again using the key identity leaves
\[
\boxed{\mathrm{Vega}=\frac{\partial C}{\partial\sigma}
=S\phi(d_1)\sqrt T.}
\]

\emph{Rho.} Since
$\partial_rd_1=\partial_rd_2=\sqrt T/\sigma$, the chain-rule terms cancel and
the derivative of the discount factor remains:
\[
\boxed{\mathrm{Rho}=\frac{\partial C}{\partial r}
=KT e^{-rT}\Phi(d_2).}
\]

\emph{Theta.} Differentiation with respect to time to maturity yields
\[
\frac{\partial C}{\partial T}
=\frac{S\phi(d_1)\sigma}{2\sqrt T}
+rKe^{-rT}\Phi(d_2).
\]
Consequently, calendar-time Theta is
\[
\boxed{\Theta=-\frac{\partial C}{\partial T}
=-\frac{S\phi(d_1)\sigma}{2\sqrt T}
-rKe^{-rT}\Phi(d_2).}
\]
\textbf{Question 4 -- Short-maturity at-the-money approximation.}

\textbf{(a)} At the money forward means $K=Se^{rT}$. Discounting both sides
shows that this is equivalent to $Ke^{-rT}=S$.

\textbf{(b)} Substitution in the definitions of $d_1$ and $d_2$ gives
\[
d_1=\frac{\sigma\sqrt T}{2},\qquad
d_2=-\frac{\sigma\sqrt T}{2}.
\]

\textbf{(c)} Since $Ke^{-rT}=S$, the call price becomes
\[
C=S\left[\Phi\left(\frac{\sigma\sqrt T}{2}\right)
-\Phi\left(-\frac{\sigma\sqrt T}{2}\right)\right].
\]
Using $\Phi(-z)=1-\Phi(z)$ and
$\Phi(z)=\frac12+z/\sqrt{2\pi}+O(z^3)$,

\textbf{(d)} With $z=\sigma\sqrt T/2$, the expansion gives
\[
C=S\left[
2\left(\frac12+\frac{\sigma\sqrt T}{2\sqrt{2\pi}}
+O(T^{3/2})\right)-1\right].
\]
Therefore
\solution{$\displaystyle
C_{\rm ATM,fwd}=S\sigma\sqrt{\frac{T}{2\pi}}+O(T^{3/2})$
as $T\downarrow0$.}

\textbf{(e) Interpretation of the $\sqrt T$ scale.} Over a short horizon $T$,
the typical random move of the underlying in a diffusion model is of size
\[
S\sigma\sqrt T.
\]
An at-the-money option has almost no intrinsic value, so its price is mainly
the value of the positive part of this uncertain short-horizon move. This is
why its leading-order time value inherits the same $\sqrt T$ scale:
\[
C_{\rm ATM,fwd}\sim
\frac{1}{\sqrt{2\pi}}S\sigma\sqrt T.
\]
The dependence on maturity is therefore not linear. Dividing $T$ by four
divides the leading-order option value by two; halving $T$ divides it by
$\sqrt2$, not by two. Equivalently, differentiating the approximation gives
\[
\frac{\partial C}{\partial T}
\sim\frac{S\sigma}{2\sqrt{2\pi T}},
\qquad
\Theta=-\frac{\partial C}{\partial T}
\sim-\frac{S\sigma}{2\sqrt{2\pi T}}.
\]
Thus the magnitude of time decay becomes very large as $T\downarrow0$ at the
money. This singular behavior is specific to the immediate neighborhood of
the strike; away from the money, the short-maturity asymptotics are different.

\section{Numerical Schemes for ODEs}

\paragraph{Exercise 9 -- Euler accuracy and the classical inequalities.}
For $y'=-2y$,
\[
y_{n+1}^{E}=(1-2h)y_n^E,\qquad
y_{n+1}^{I}=\frac{1}{1+2h}y_n^I,
\]
so
\solution{\[
y_n^E=(1-2h)^n,\qquad y_n^I=(1+2h)^{-n},\qquad y(t_n)=e^{-2nh}.
\]}
Both methods have global error $O(h)$, so halving $h$ asymptotically halves
the maximum error. Explicit Euler is stable precisely when
$|1-2h|<1$, hence $0<h<1$. The implicit multiplier satisfies
$0<(1+2h)^{-1}<1$ for every $h>0$. Stability controls growth, not the size of
the discretization error; a large stable step can therefore be inaccurate.

\section{Solving Scalar and Multidimensional ODEs}

\paragraph{Exercise 10 -- Scalar linear equations.}
Applying variation of constants or the integrating-factor method gives
\solution{\[
y_1(t)=\frac{t}{2}-\frac18-\frac78e^{-4t},
\qquad
y_2(t)=(1+t)\left(2+t+\frac{t^2}{2}\right),
\qquad
y_3(t)=e^{2t}-e^t.
\]}
For the second equation, the integrating factor is $(1+t)^{-1}$ and
$[y/(1+t)]'=1+t$. Direct differentiation verifies all three initial values
and equations (the second solution is considered on the interval $t>-1$).

\paragraph{Exercise 11 -- Nonlinear scalar equations.}
Separation of variables gives
\solution{\[
y'=y(1-y):\quad
y(t)=\frac{1}{1+\frac{1-y_0}{y_0}e^{-t}},\qquad t\in\R.
\]}
The equilibria are $0$ (unstable) and $1$ (asymptotically stable). Similarly,
\[
y'=-y^3:\quad y(t)=\frac{1}{\sqrt{1+2t}},
\qquad t\in\left(-\frac12,\infty\right).
\]
Its only equilibrium is $0$, which is asymptotically stable, although the
linearization vanishes. Finally,
\[
y'=t(1+y^2):\quad \arctan y=\frac{t^2}{2},\qquad
y(t)=\tan\left(\frac{t^2}{2}\right).
\]
The maximal interval containing zero is $(-\sqrt\pi,\sqrt\pi)$. The last
equation is nonautonomous, so a one-dimensional autonomous phase line is not
applicable to it.

\paragraph{Exercise 12 -- Linear systems in dimension two.}
The eigenpairs may be chosen as
\[
\lambda_1=-2,\ v_1=(1,0)^\top,\qquad
\lambda_2=-4,\ v_2=(1,-2)^\top.
\]
Since $x(0)=2v_1-v_2$,
\solution{\[
x(t)=2e^{-2t}v_1-e^{-4t}v_2
=\binom{2e^{-2t}-e^{-4t}}{2e^{-4t}}.
\]}
This is exactly $e^{tA}x(0)$. The numerical updates are
\[
x_{n+1}^{E}=(I+hA)x_n^E,
\qquad x_{n+1}^{I}=(I-hA)^{-1}x_n^I.
\]
Their errors against the displayed exact solution decay linearly with $h$.

\paragraph{Exercise 13 -- Forced multidimensional system.}
Solving $Ax^\star+b=0$ gives $x^\star=(1,1)^\top$. With
$z=x-x^\star$, $z'=Az$ and $z(0)=(-1,1)^\top$. Componentwise,
$z_1=-e^{-2t}$ and $z_2=e^{-2t}$, hence
\solution{\[
x(t)=\binom{1-e^{-2t}}{1+e^{-2t}}\longrightarrow\binom11=x^\star.
\]}
Explicit Euler uses $x_{n+1}=(I+hA)x_n+hb$; implicit Euler uses
$x_{n+1}=(I-hA)^{-1}(x_n+hb)$. Both converge to the same equilibrium when
their stability conditions are respected.

\paragraph{Exercise 14 -- Second-order equation as a system.}
The characteristic polynomial is $(r+1)(r+3)$, so the initial data yield
\solution{$\displaystyle y(t)=e^{-t}-e^{-3t}$.}
With $z=(y,y')^\top$,
\[
z'=\begin{pmatrix}0&1\\-3&-4\end{pmatrix}z,
\qquad z(0)=\binom02.
\]
The companion matrix has the same eigenvalues $-1$ and $-3$. Explicit Euler
is $z_{n+1}=(I+hA_c)z_n$; comparing its first component with
$e^{-t_n}-e^{-3t_n}$ verifies first-order convergence.

\end{document}
